\clearpage
\begingroup
\fancyhead[L]{\small NACHLASS UND BRIEFWECHSEL}
\renewcommand{\headrulewidth}{0pt}
\clearpage
\SourceLeaf{584}
\pdfbookmark[1]{Nachlass VIII. Störungen der Pallas durch Saturn}{pallas-nachlass-VIII}
\begin{center}
\large NACHLASS UND BRIEFWECHSEL.
\par\bigskip
\Large [VIII. STÖRUNGEN DER PALLAS\par DURCH SATURN.]
\par\bigskip
\normalsize [1.]
\end{center}
\pdfbookmark[2]{[1]}{pallas-VIII-fragment1}
\begingroup
\small
\makeatletter\tagsleft@true\let\veqno\@@leqno\makeatother
\par\noindent\hspace*{\parindent}[Die Störungen durch Saturn wurden von \textsc{Nicolai} im Wesentlichen in derselben Weise
gerechnet, wie die durch Jupiter, indem die störenden Kräfte nach der interpolatorischen
Methode nach Reihen mit den Argumenten \(M\) und \(M'\) entwickelt wurden. Indessen wurden nicht
die drei Grössen \(T,V,W\) als Componenten der störenden Kräfte verwandt, sondern diese
wurden in drei rechtwinklige Componenten \(X,Y,Z\) zerlegt, von denen die erste in die
Richtung der Apsidenlinie der Pallas fällt, die zweite senkrecht dazu in der Bahnebene der
Pallas liegt und die dritte senkrecht zu dieser Bahnebene steht. Sind \(x,y,z,x',y',z'\)
die rechtwinkligen Coordinaten beider Planeten in bezug auf die eben genannten Axen, so sind
diese Componenten
\[
\tag*{1)}
\begin{aligned}
X&=\mu\left(\frac{x'}{r'^3}-\frac{x'-x}{\rho^3}\right)\\[.6em]
Y&=\mu\left(\frac{y'}{r'^3}-\frac{y'-y}{\rho^3}\right)\\[.6em]
Z&=\mu\left(\frac{z'}{r'^3}-\frac{z'-z}{\rho^3}\right)\qquad(z=0)
\end{aligned}
\]
und ihre Beziehungen zu den früher benutzten sind:
\[
\begin{aligned}
\frac{ar}{\cos\varphi}\bigl(Y\cos(v-\tilde\omega)-X\sin(v-\tilde\omega)\bigr)&=V\\[.6em]
\frac{a}{\cos\varphi}\bigl(Y\sin(v-\tilde\omega)+X\cos(v-\tilde\omega)\bigr)&=T\\[.6em]
-\frac{ar}{\cos\varphi}Z&=W.
\end{aligned}
\]
\par\noindent\hspace*{\parindent}Die Störungen der Elemente bestimmen sich dann aus den Gleichungen
\SourceLeaf{585}
\[
\tag*{2)}
\begin{aligned}
\frac{\mathrm{d}i}{\mathrm{d}M}
&=-\frac{ar\cos(v-\AscendingNode)}{\cos\varphi}Z\\[.6em]
\frac{\mathrm{d}\AscendingNode}{\mathrm{d}M}
&=-\frac{ar\sin(v-\AscendingNode)}{\cos\varphi\sin i}Z\\[.6em]
\frac{\mathrm{d}\varphi}{\mathrm{d}M}
&=aa\sin(v-\tilde\omega)\mathord{.}\cos E\mathord{.}X
 -aa\bigl(1+\cos(v-\tilde\omega)\mathord{.}\cos E\bigr)Y\\[.6em]
\frac{\mathrm{d}\tilde\omega}{\mathrm{d}M}
&=\frac{aa}{e}\bigl(\cos\varphi+\sin(v-\tilde\omega)\mathord{.}\sin E\bigr)X
 -\frac{aa\cos(v-\tilde\omega)\sin E}{e}Y+A\\[.6em]
\frac{\mathrm{d}n}{\mathrm{d}M}
&=\frac{3a^3n\cos\varphi\cos E}{r}Y-\frac{3a^3n\sin E}{r}X\\[.6em]
\frac{\mathrm{d}\varepsilon}{\mathrm{d}M}
&=\frac{\mathrm{d}\tilde\omega}{\mathrm{d}M}
 -\frac{rr\left(\dfrac{\mathrm{d}\tilde\omega}{\mathrm{d}M}-A\right)}{aa\cos\varphi}
 -\frac{\sin E(2-ee-e\cos E)}{\cos\varphi}\frac{\mathrm{d}\varphi}{\mathrm{d}M}\\[.6em]
A&=-\frac{ar\,\mathrm{tg}\,\tfrac{1}{2}i\sin(v-\AscendingNode)}{\cos\varphi}Z,
\end{aligned}
\]
welche sich leicht aus den Artikeln 8--12 der Exposition d'une nouvelle méthode, S. 450--455
dieses Bandes, ableiten lassen, und wo die Bezeichnungen die S. 473--474 dieses Bandes
erläuterte Bedeutung haben.]
\par\endgroup
\clearpage
\endgroup
